\documentclass[10pt,a4paper]{article}

\usepackage[utf8]{inputenc}
\usepackage[T1]{fontenc}
\usepackage[ngerman]{babel}
\usepackage[margin=2cm]{geometry}
\usepackage{graphicx}
\usepackage{hyperref}

% Platzhalter-Box fuer eine Abbildung, die noch von Hand erstellt werden muss.
% #1 = Bildunterschrift; der Inhalt beschreibt, was hineingehoert.
\newsavebox{\placebox}
\newenvironment{placeholder}[1]
  {\par\medskip\noindent\begin{center}
   \setlength{\fboxsep}{6pt}%
   \begin{lrbox}{\placebox}%
   \begin{minipage}{0.88\linewidth}
   \centering\textbf{[ PLATZHALTER: #1 ]}\\[0.8em]
   \raggedright\footnotesize}
  {\end{minipage}\end{lrbox}\fbox{\usebox{\placebox}}\end{center}\medskip}

\title{\textbf{mapping\_pathfinding}\\[0.3em]
       \large Kartierung, Lokalisierung und Pfadplanung f\"ur die Duckiebot-City-Challenge}
\author{Nico Sander}
\date{\today}
\usepackage{titlesec}
\titlespacing*{\section}{0pt}{1.2ex plus .5ex}{0.8ex}
\titlespacing*{\subsection}{0pt}{1ex plus .3ex}{0.5ex}
\setlength{\parskip}{0.4ex}

\begin{document}
\maketitle

\section{Systemüberblick}

\texttt{mapping\_pathfinding} ist ein ROS-1-Paket (Noetic), mit dem ein
Duckiebot eine kleine Stadt aus Straßen und Kreuzungen selbstständig erkundet,
eine Karte davon aufbaut und anschließend eine angesagte Folge von Toren in der
schnellsten gefundenen Reihenfolge abfährt.

Die Mission besteht aus zwei Phasen:

\begin{description}
  \item[MAPPING (ungestoppt)] Der Bot fährt jede Straße der Stadt mindestens
    einmal ab. Dabei erfasst er, welches \emph{Tor} (ein AprilTag über einer
    Straße) zu welcher Straße gehört, und misst, wie lange er für jede Straße
    und jedes Abbiegemanöver tatsächlich braucht.
  \item[GATE\_RUN (gestoppt)] Die Torreihenfolge wird vor Ort als Liste von
    Tag-IDs angesagt. Der Bot wird an einer bekannten Straße abgesetzt, plant
    anhand der Karte und der gemessenen Zeiten aus der Mapping-Phase eine Route
    durch die Tore in dieser Reihenfolge und fährt sie ab.
  \item[DONE] Nach dem letzten Tor hält der Bot an und meldet die benötigte
    Zeit.
\end{description}

Die Stadt wird als Graph modelliert: Knoten sind Kreuzungen, Kanten sind
Straßen. Jede Kreuzung hat bis zu vier nummerierte \emph{Ports} (1--4), die
zyklisch angeordnet sind -- für einen Bot, der über Port $p$ einfährt, liegt
Port $p+1$ rechts und $p-1$ links. Da Wendemanöver verboten sind, wird nicht
über Knoten geplant, sondern über gerichtete Zustände \texttt{(Knoten,
Einfahrt-Port)} -- ``ich nähere mich dieser Kreuzung und bin über diesen Port
eingefahren''. Dieser Zustand identifiziert zugleich die gerade befahrene
Straße, was es dem Planer erlaubt, einzelne Straßen und Tore gezielt
anzusteuern. Die Karte selbst liegt in einer einzigen JSON-Datei
(\texttt{config/city.json}), die alle Nodes lesen -- es gibt also immer nur
genau eine Definition der Stadt.

Das Paket trennt bewusst zwischen dünnen ROS-Nodes und einem ROS-freien Kern:
sämtliche Graph-, Planungs-, Zeitmess- und Filterlogik steckt in reinen
Python-Modulen ohne ROS-Import. Dadurch ist der gesamte Entscheidungskern durch
Offline-Unittests abgedeckt (derzeit 209), die ohne Roboter laufen.

\begin{placeholder}{ROS-Node-Graph}
  Ein Blockdiagramm des laufenden ROS-Graphen. Pro Node ein Kasten
  (\texttt{detect\_lane}, \texttt{detect\_signs}, \texttt{detect\_intersection},
  \texttt{switch\_control}, \texttt{control\_wheels},
  \texttt{mapping\_pathplanning}, \texttt{visualization}, \texttt{dashboard}),
  dazu Kamera und Radansteuerung als externe Kästen, verbunden durch Pfeile mit
  dem jeweiligen Topic-Namen. Die wichtigen Kanten sind:
  Kamera $\rightarrow$ \texttt{detect\_lane} und \texttt{detect\_signs};
  \texttt{/detect/lane\_borders} $\rightarrow$ \texttt{detect\_intersection};
  \texttt{/detect/intersection} und \texttt{/detect/sign} $\rightarrow$
  \texttt{switch\_control}; \texttt{/switch/mode} und
  \texttt{/switch/turn\_direction} $\rightarrow$ \texttt{control\_wheels}
  \emph{und} $\rightarrow$ \texttt{mapping\_pathplanning};
  \texttt{/detect/sign\_detections} $\rightarrow$
  \texttt{mapping\_pathplanning}; und die Kante, die den Regelkreis schließt:
  \texttt{/plan/turn\_command} $\rightarrow$ \texttt{switch\_control}.
  Ebenfalls einzeichnen: \texttt{/mapping/state} $\rightarrow$
  \texttt{visualization}.
  Optisch hervorheben lohnt sich die Schichtung -- Wahrnehmung (links),
  Verhalten (Mitte), Mission/Planung (rechts) -- und die Tatsache, dass die
  Missionsschicht auf die Verhaltensschicht zurückwirkt. Ein Screenshot von
  \texttt{rqt\_graph} tut es zur Not auch, eine handgezeichnete Version ist aber
  deutlich lesbarer, weil sie die Debug-Topics weglassen kann.
\end{placeholder}

\section{Beschreibung der Kernkomponenten}

\subsection{Wahrnehmung}

\paragraph{\texttt{detect\_lane.py}} Segmentiert das Kamerabild mit einem
PyTorch-U-Net in weiße Fahrbahnmarkierungen, gelbe Mittellinie und rote
Haltelinien. Daraus berechnet der Node den Querabstandsfehler -- wie weit der
Bot aus der Spurmitte steht -- und veröffentlicht ihn auf
\texttt{/detect/lane}. Zusätzlich publiziert er die geometrischen Spurgrenzen
(inklusive Position einer roten Linie im Bild) auf
\texttt{/detect/lane\_borders} sowie Debug-Bilder für das Dashboard. Die
gesamte Wahrnehmungspipeline läuft mit 30\,Hz auf dem Laptop, etwa 20\,ms pro
Bild bei einem Budget von 33\,ms.

\paragraph{\texttt{detect\_signs.py}} Erkennt AprilTags mit
\texttt{pupil\_apriltags} und publiziert auf zwei Topics, weil die beiden
Abnehmer Unterschiedliches brauchen. \texttt{/detect/sign} enthält nur die ID
des nächstgelegenen Tags und ist ungefiltert -- \texttt{switch\_control} muss
Kreuzungsschilder in jeder Größe sehen. \texttt{/detect/sign\_detections}
enthält \emph{jede} Detektion mit Eckpunkten, Fläche in Pixeln und einer
Annahme-/Verwerfen-Entscheidung; der Missions-Node braucht die Fläche, um zu
unterscheiden, ob ein Tor auf der aktuellen Straße hängt oder erst hinter der
Kreuzung.

\paragraph{\texttt{detect\_intersection.py}} Ein kleiner Logik-Node, der aus der
Spurgeometrie einen Kreuzungszustand ableitet: keine Kreuzung, Kreuzung nähert
sich, Haltelinie erreicht. Das geschieht allein über die vertikale Position der
erkannten roten Linie, ohne das Bild erneut anzufassen.

\subsection{Verhalten und Aktorik}

\paragraph{\texttt{switch\_control.py}} Der Verhaltensautomat. Er entscheidet,
\emph{was der Bot gerade tut} -- Spur folgen, sich einer Kreuzung nähern, an der
Haltelinie stehen oder die Kreuzung queren -- und veröffentlicht das als
\texttt{/switch/mode} zusammen mit der Abbiegerichtung auf
\texttt{/switch/turn\_direction}. Die Richtung kommt im Normalfall vom
Missions-Node über \texttt{/plan/turn\_command}; ist dieses Kommando veraltet
oder fehlt es, fällt der Node auf sein älteres Verhalten zurück und wählt die
Richtung anhand des Kreuzungsschilds -- ein ausgefallener Missions-Node führt so
zu degradiertem statt zu gar keinem Fahrverhalten. Kreuzungsquerungen sind
geschlossen geregelt: das Manöver endet, sobald wieder Fahrbahnmarkierungen zu
sehen sind, nicht wenn ein Timer abläuft -- die konfigurierte Dauer ist nur ein
Timeout. Genau das macht den Stack tolerant dagegen, dass der Bot nicht exakt
rechtwinklig an der Haltelinie steht.

\paragraph{\texttt{control\_wheels.py}} Der einzige Node, der mit den Motoren
spricht. Er übersetzt den aktuellen Fahrmodus in Radkommandos: beim Spurfolgen
läuft ein PID-Regler auf dem Querabstandsfehler (mit tiefpassgefiltertem
D-Anteil, damit er bei 30\,Hz stabil bleibt), während einer Querung fährt er
einen festen Bogen für die kommandierte Richtung, bis \texttt{switch\_control}
das Manöver beendet. Gemessene Latenz von Bildaufnahme bis Kommando: rund
75\,ms.

\subsection{Mission, Kartierung und Planung}

\paragraph{\texttt{mapping\_pathplanning.py}} Der Node, dem die Mission gehört,
und die größte Einzelkomponente des Pakets. Er erledigt vier Dinge. \emph{(1)
Lokalisierung:} Koppelnavigation auf dem Graphen -- die geglaubte Position
rückt um eine Straße weiter, sobald \texttt{switch\_control} eine
abgeschlossene Querung meldet. \emph{(2) Kartierung:} er schreibt Tore auf die
Straßen, auf denen sie gesehen wurden. \emph{(3) Planung:} nach jeder Querung
wird die gesamte Restroute neu geplant -- bei dieser Stadtgröße billig, und es
bedeutet, dass eine korrigierte Position sofort zu einer korrigierten Route
führt. \emph{(4) Kommandieren:} er publiziert das nächste Abbiegekommando auf
\texttt{/plan/turn\_command} und schließt damit den Regelkreis. Zusätzlich
veröffentlicht er den vollständigen Missionszustand (Karte, Tore, Pfad,
Position, Phase) als JSON auf \texttt{/mapping/state} und bietet einen Service
\texttt{start\_gate\_run} für den Phasenwechsel an.

\paragraph{\texttt{planner.py} (ROS-frei)} Dijkstra über die gerichteten
Zustände \texttt{(Knoten, Einfahrt-Port)}. Wendemanöver werden nicht durch eine
Prüfung verboten, sondern sind strukturell nicht ausdrückbar, weil keine
Richtungstabelle einen Port auf sich selbst abbildet. Kanten sind mit einer
geschätzten Dauer gewichtet -- ein Zug kostet \texttt{Halt + Abbiegen +
Straße} -- damit die Suche wirklich schnellere und nicht bloß kürzere Routen
bevorzugt; eine Rechtskurve ist deutlich billiger als eine Linkskurve. In der
Mapping-Phase treibt dasselbe Modul eine gierige Erkundung nach der jeweils
nächstgelegenen noch nicht befahrenen Straße.

\paragraph{\texttt{city\_map.py} (ROS-frei)} Der Graph selbst: Port-Geometrie,
Laden und Validieren von \texttt{city.json} sowie \texttt{GraphMap}, der
Live-Kartenzustand, der festhält, welche Straßen befahren wurden und welches Tor
auf welcher Straße hängt. Kantenschlüssel sind richtungsunabhängig
(\texttt{A1\_\_B1} von beiden Enden), womit sich ``ein Tor gilt für die ganze
Straße'' von selbst ergibt.

\paragraph{\texttt{gate\_detection.py} (ROS-frei)} Entscheidet, welchen
Tor-Sichtungen geglaubt werden darf. Ein Tor wird dauerhaft auf eine Straße
geschrieben -- der erste Eintrag gewinnt -- weil die Kamera über die Kreuzung
hinaus sieht und spätere Sichtungen aus größerer Entfernung und schlechterem
Winkel stammen. Drei Filter stehen vor diesem unwiderruflichen Schreibvorgang:
im Stand und während einer Querung wird gar nicht kartiert; die Fläche des Tags
muss über einem aus aufgezeichneten Fahrten kalibrierten Schwellwert liegen
(\texttt{gate\_min\_area}, derzeit 900\,px$^2$); und dasselbe Tag muss über
mehrere aufeinanderfolgende Bilder auf derselben Straße gesehen werden.

\paragraph{\texttt{run\_timing.py} (ROS-frei)} Misst, wie lange jede Straße und
jedes Abbiegemanöver tatsächlich dauert. Als gemeinsame Zeitbasis dienen die von
\texttt{switch\_control} publizierten Fahrmodus-Übergänge, sodass sich die
gemessenen Teile lückenlos und überlappungsfrei zu einem vollen Zug
zusammensetzen. Die Mapping-Phase ist ungestoppt, dort zu messen kostet also
nichts; der gestoppte Lauf wird anschließend gegen reale Zahlen statt gegen eine
pauschale Schätzung geplant. Mehrfachmessungen einer Straße werden über den
Median zusammengefasst.

\paragraph{\texttt{crossing.py} und \texttt{graph\_layout.py} (ROS-frei)} Zwei
kleine Hilfsmodule: das erste entscheidet, wann eine Kreuzungsquerung beendet
ist (Markierungen wieder sichtbar, nach einer minimalen Blindzeit), das zweite
berechnet die Zeichengeometrie der Karte -- mit der getesteten Eigenschaft, dass
sich gezeichnete Straßen außerhalb eines Knotens niemals kreuzen dürfen.

\subsection{Bedienung und Visualisierung}

\paragraph{\texttt{visualization.py}} Ein Live-Fenster, das die drei von der
Challenge geforderten Dinge sichtbar macht: die aufgebaute Karte (Straßen mit
den darauf gefundenen Toren), den geplanten Pfad (gestrichelt, mit
Schrittnummern) und die geglaubte Position des Bots (aktuelle Straße
hervorgehoben, Ring an der nächsten Kreuzung).

\paragraph{\texttt{dashboard.py}} Ein Debug-Fenster für die Wahrnehmung, das
Rohbild, Segmentierungsmasken und die AprilTag-Rahmen samt gemessener Fläche in
einer einzigen Ansicht kachelt.

\paragraph{\texttt{mission\_tui.py} / \texttt{mission\_setup.py}} Ein
menügeführtes Frontend für die gesamte Session. Der strukturell wichtige Punkt:
\emph{ein} \texttt{roslaunch} bleibt über beide Phasen hinweg am Leben -- Karte,
Tore und gemessene Zeiten liegen im Missions-Node, ein Neustart zwischen den
Phasen würde sie verwerfen. Die TUI besitzt daher den Launch-Prozess und
wechselt die Phase über Parameter und den Service-Aufruf.
\texttt{mission\_setup.py} validiert die eingegebenen Antworten (vor allem die
Startposition, die vor dem Start gegen den Stadtgraphen geprüft wird) und baut
das Launch-Kommando zusammen.

\section{Zustands- und Ablaufdiagramme}

\begin{placeholder}{Zustandsdiagramm der Missionsphasen}
  Der übergeordnete Missionsautomat von \texttt{mapping\_pathplanning.py} mit
  drei Zuständen: \textbf{MAPPING}, \textbf{GATE\_RUN} und \textbf{DONE}.
  Einzuzeichnende Übergänge: MAPPING $\rightarrow$ GATE\_RUN, beschriftet mit
  ``alle Straßen befahren \emph{und} alle Tore gesehen $\rightarrow$ Halt an der
  Haltelinie, warten auf den Service \texttt{start\_gate\_run}''; GATE\_RUN
  $\rightarrow$ DONE, beschriftet mit ``letztes Tor der angesagten Reihenfolge
  passiert''. Deutlich markieren, dass zwischen MAPPING und GATE\_RUN ein
  \emph{harter Stopp} liegt, während dessen der Bot physisch umgesetzt wird, und
  dass der Übergang aus MAPPING verlangt, dass die Tore tatsächlich
  \emph{gesehen} wurden -- nicht nur, dass die Straßen befahren wurden.
\end{placeholder}

\begin{placeholder}{Verhaltensautomat (\texttt{switch\_control})}
  Der Fahrmodus-Automat: \textbf{LANE\_FOLLOWING} $\rightarrow$
  \textbf{APPROACHING} $\rightarrow$ \textbf{STOPPED} $\rightarrow$
  \textbf{CROSSING} $\rightarrow$ zurück zu LANE\_FOLLOWING. Jeden Übergang mit
  seinem Auslöser beschriften: rote Linie oberhalb des Schwellwerts erkannt;
  Haltelinie erreicht; Haltedauer abgelaufen \emph{und} Abbiegekommando
  vorhanden; und -- der wichtigste -- die Querung endet, sobald \emph{wieder
  Fahrbahnmarkierungen sichtbar} sind, wobei die konfigurierte Abbiegedauer nur
  als Timeout wirkt. Zusätzlich das \texttt{HALT}-Kommando als Übergang in einen
  Wartezustand einzeichnen sowie \texttt{/plan/resume\_driving} als jene Kante,
  die einen umgesetzten Bot ohne Querungsmanöver zurück in LANE\_FOLLOWING
  bringt. Es lohnt sich, im selben Diagramm zu annotieren, welches Zeitintervall
  \texttt{run\_timing.py} als ``Abbiegezeit'' und welches als ``Straßenzeit''
  misst.
\end{placeholder}

\begin{placeholder}{Ablaufdiagramm der Tor-Sichtungsentscheidung}
  Ein Flussdiagramm für eine einzelne eingehende Detektion aus
  \texttt{/detect/sign\_detections}, das entweder in ``Tor auf Straße
  schreiben'' oder in ``verwerfen'' mündet. Entscheidungsrauten der Reihe nach:
  Fährt der Bot überhaupt (nicht gestoppt, nicht querend)? Ist die Tag-ID
  überhaupt ein Tor? Ist ihre Fläche $\geq$ \texttt{gate\_min\_area}? Wurde
  dasselbe Paar (Straße, Tor) über \texttt{gate\_confirm\_frames}
  aufeinanderfolgende Detektionen gesehen? Trägt diese Straße bereits ein Tor
  (falls ja: verwerfen -- der erste Eintrag ist endgültig)? Dieses Diagramm ist
  die beste Stelle, um zu zeigen, \emph{warum} es jeden Filter gibt; eine
  einzeilige Anmerkung pro Raute genügt.
\end{placeholder}

\begin{placeholder}{Regelkreis / Sequenz an einer Kreuzung}
  Ein Sequenzdiagramm über eine Kreuzung hinweg, mit den Spuren
  Kamera $\rightarrow$ \texttt{detect\_*} $\rightarrow$ \texttt{switch\_control}
  $\rightarrow$ \texttt{control\_wheels} und \texttt{mapping\_pathplanning} als
  paralleler Spur. Der Ablauf: rote Linie erkannt, Halt, Missions-Node plant neu
  und publiziert das Abbiegekommando, \texttt{switch\_control} geht mit dieser
  Richtung in CROSSING, \texttt{control\_wheels} fährt den Bogen,
  Fahrbahnmarkierungen tauchen wieder auf, \texttt{switch\_control} meldet die
  abgeschlossene Querung, der Missions-Node rückt seine geglaubte Position um
  eine Straße weiter und plant neu. Dieses Diagramm macht den geschlossenen
  Regelkreis lesbar -- und die 30-Hz-Innenschleife innerhalb der äußeren
  Schleife pro Kreuzung.
\end{placeholder}

\section{Konfiguration und Tests}

Alle abgestimmten Konstanten liegen in \texttt{config/config.json} (Zeiten,
PID-Parameter, Schwellwerte, Kostenmodell des Planers); die Karte ist
\texttt{config/city.json}, die rein kosmetische Zuordnung Tor-ID $\rightarrow$
Farbe steht in \texttt{config/gates.json}. Launch-Argumente überschreiben
ausgewählte Werte für einen einzelnen Lauf.

Da der Entscheidungskern ROS-frei ist, wird er offline verifiziert:
Port-Geometrie und die Invariante ``kein Wendemanöver'', Validierung der Stadt,
die Filter für Tor-Sichtungen, Erreichbarkeit im Planer aus jedem Startzustand,
alle Torreihenfolgen aus allen Startzuständen, gierige Abdeckung, vollständige
Missionssimulationen und synthetische Gitterstädte. Zusätzlich betreibt
\texttt{launch/simulate.launch} den echten Missions-Node gegen einen simulierten
Fahrer und testet damit die tatsächliche ROS-Schicht -- Topics, Timer, den
Phasenwechsel -- ohne Kamera, Motoren oder Roboter.

\end{document}
